\documentclass[11pt,a4paper]{article}
\usepackage[T1]{fontenc}
\usepackage[utf8]{inputenc}
\usepackage{lmodern}
\usepackage{amsmath,amssymb,amsthm,mathtools,mathrsfs,bm}
\usepackage[margin=1in]{geometry}
\usepackage{microtype,booktabs,array}
\usepackage[hidelinks]{hyperref}
\usepackage[nameinlink,noabbrev]{cleveref}
\numberwithin{equation}{section}
\newtheorem{theorem}{Theorem}[section]
\newtheorem{lemma}[theorem]{Lemma}
\newtheorem{proposition}[theorem]{Proposition}
\newtheorem{corollary}[theorem]{Corollary}
\theoremstyle{definition}
\newtheorem{definition}[theorem]{Definition}
\newtheorem{remark}[theorem]{Remark}
\allowdisplaybreaks[1]
\title{Non-injectivity of a Terminal-Temperature Map in a Coupled Heat--Memory Model}
\author{Vipin Gupta\thanks{Affiliation to be confirmed.}}
\date{}
\begin{document}
\maketitle
\begin{abstract}
For a frozen linear heat--memory system with zero initial data and antisymmetric boundary control, zero input trivially gives zero terminal temperature. We study instead non-injectivity: an infinite-dimensional space of smooth compatible controls has zero terminal temperature, including all its nonzero members. The argument uses a classical Fattorini--Russell biorthogonal estimate, M\"untz incompleteness, and a Fredholm range intersection for an exact coupled moment operator. The terminal memory values obey an entire-function trace constraint; we give necessary asymptotic conditions and control-norm lower bounds, but do not determine the exact memory range or its closure. Simultaneous vanishing of temperature and memory forces zero input by the accumulating-zero mechanism of Ivanov and Pandolfi. This is not a nonlinear preparation theorem or a controllability theorem for arbitrary initial states. Independent review and a complete novelty comparison remain outstanding.
\end{abstract}
\noindent\textbf{Keywords:} exponential moments; biorthogonal families; M\"untz theorem; compact operators; heat equations with memory.\\
\textbf{MSC 2020:} 44A10, 30E05, 93B05, 35R09.
\section{Introduction and source comparison}\label{sec:intro}
With zero initial data the zero control makes the terminal temperature zero. The point considered here is instead the existence of many \emph{nonzero} controls with the same zero temperature. This is a non-injectivity question for an output map, not null controllability from arbitrary initial states and not full-state return to rest. The controls considered below belong to the stated smooth compatible class.

Fattorini and Russell's 1971 paper \cite{FattoriniRussell1971} uses exponential moments for a one-dimensional scalar parabolic equation; its Theorems 3.3--3.4 treat reachable and null-controllable states including smooth compactly supported controls under the stated decay hypotheses. Their 1974 Theorems 1.2 and 1.5(b) \cite{FattoriniRussell1974} give the fixed-time negative-exponential biorthogonal estimate used here. Neither source directly addresses the coupled slow branch. The 2011 theorem of Tenenbaum and Tucsnak \cite{TenenbaumTucsnak2011} gives arbitrary-initial-data null controllability and cost estimates under a self-adjoint spectral hypothesis and a spectral inequality. Its eigenvalues tend to infinity; this does not cover our full accumulating spectrum directly.

Ivanov and Pandolfi \cite{IvanovPandolfi2009} study rest from arbitrary initial states for a heat equation in which the memory acts on the spatial Laplacian. Their Obstruction 1 uses accumulation of characteristic roots and analyticity to obstruct interpolation. We credit that mechanism when establishing that zero temperature and zero memory together force zero control from zero start; the conclusion is not claimed as new. Our bounded ODE coupling yields a different equation and target, and their negative rest theorem is not invoked as a theorem about our model. Their paper also gives conditional negative results through a completeness obstruction and fractional-kernel arguments.

Guerrero and Imanuvilov \cite{Guerrero2013} show non-null-controllability for some nonzero initial data in their specified heat--memory equations, including a model with memory acting on the state itself, using adjoint observability. Halanay and Pandolfi \cite{HalanayPandolfi2015}, for the Coleman--Gurtin equation and their stated control and terminal-trace hypotheses, obtain approximate temperature controllability and failure of temperature null controllability for some initial data when the relevant resolvent is nonzero. Neither statement settles the exact zero-start memory image below. Ammar Khodja et al.~\cite{Ammar2014} relate minimal time to condensation for a Riesz-basis spectral problem under sector and reciprocal-summability assumptions; the bounded accumulating branch here violates the latter for the combined spectrum. Shklyar's preprint \cite{Shklyar2024} studies partial null controllability by moments under a no-finite-accumulation spectrum assumption. Neither result directly settles our problem. These targeted comparisons do not prove originality or exhaust the literature.

The following sections formulate the frozen model, establish the paired moments, and give an abstract kernel construction. The memory trace restrictions are necessary conditions, not an exact characterization or a density theorem. The construction and its physical interpretation remain subject to independent review.

\section{Governing model and preliminaries}\label{sec:model}
For the interpretation, consider real fields on \(0<x<L\), \(-T<s<0\), satisfying the frozen equations
\begin{equation}\label{eq:pde}
M Y_s=\kappa Y_{xx}+B_0Y+SP,\qquad P_s=-\mu P+\gamma Y,
\end{equation}
with \(M,\kappa,L,\mu>0\), \(d=\gamma S/M>0\), and real constant \(B_0,S,\gamma\). Initially \(Y(\cdot,-T)=P(\cdot,-T)=0\); the prescribed boundary values are \(Y(0,s)=a(s)\), \(Y(L,s)=-a(s)\). We work in the antisymmetric subspace about \(x=L/2\). This is a linear frozen model, not an assertion about nonlinear reachability or prescribed physical parameter values.

For the interpretation below take a smooth control with \(a(-T)=a(0)=0\), and use a smooth compatible boundary lift \(Y=Z+a(s)(1-2x/L)\). The homogeneous-Dirichlet heat generator on \(L^2(0,L)\), coupled to the bounded memory ODE, has a bounded off-diagonal perturbation. The lift supplies continuous \(L^2\)-valued forcing, so standard semigroup variation of constants gives a continuous \(L^2\)-valued pair. The basis for the odd subspace is
\begin{equation}\label{eq:basis}
\phi_n(x)=\sqrt{2/L}\sin(q_nx),\qquad q_n=2\pi n/L,\quad n\ge1.
\end{equation}
All terminal identities below are understood in this solution class, not as a general endpoint trace theorem for arbitrary \(L^2\) boundary data. The physical relation \(\mu B_0+\gamma S=0\), when imposed, is an additional condition; it is not assumed for the abstract theorem.

\section{Modal decomposition and exact moments}\label{sec:modal}
Set \(Y_n=\langle Y,\phi_n\rangle\), \(P_n=\langle P,\phi_n\rangle\). Integration by parts and \(\phi_n'(0)=\phi_n'(L)=\sqrt{2/L}q_n\) yield
\begin{equation}\label{eq:ibp}
\int_0^L Y_{xx}\phi_n\,dx=-q_n^2Y_n+2a(s)\sqrt{2/L}\,q_n.
\end{equation}
Consequently, with \(A_n=(\kappa q_n^2-B_0)/M\) and \(C_n=2\kappa\sqrt{2/L}\,q_n/M\),
\begin{equation}\label{eq:modalode}
\frac d{ds}\binom{Y_n}{P_n}=D_n\binom{Y_n}{P_n}+\binom{C_na}{0},\qquad
D_n=\begin{pmatrix}-A_n&S/M\\\gamma&-\mu\end{pmatrix}.
\end{equation}
The characteristic equation is \((\lambda+A_n)(\lambda+\mu)-d=0\). Define
\begin{equation}\label{eq:roots}
b_n=\frac{A_n-\mu+\sqrt{(A_n-\mu)^2+4d}}2>0,\qquad \delta_n=\frac d{b_n}>0.
\end{equation}
Then the roots are \(-\mu-b_n\), \(-\mu+\delta_n\), and \(b_n-\delta_n=A_n-\mu\). The root \(b_n\) increases strictly with \(A_n\). Writing \(\alpha=\kappa(2\pi/L)^2/M>0\), \(\beta=-B_0/M-\mu\), gives \(b_n=\alpha n^2+\beta+O(n^{-2})\); its square-root gaps tend to \(\sqrt\alpha\). Strict increase and this limit give a positive global gap. If \(\mu B_0+\gamma S=0\), the determinant \(\mu A_n-d=\mu\kappa q_n^2/M\) and the trace are positive and negative respectively, so both roots are negative.

For \(f(t)=a(-t)\) and \(F(z)=\int_0^T e^{zt}f(t)dt\), variation of constants gives \(\binom{Y_n(0)}{P_n(0)}=C_n\int_0^T e^{tD_n}\binom10 f(t)dt\). The two spectral projectors give
\begin{equation}\label{eq:projector}
e^{tD_n}\binom10=\frac1{b_n+\delta_n}
\binom{\delta_ne^{(-\mu+\delta_n)t}+b_ne^{(-\mu-b_n)t}}
{\gamma(e^{(-\mu+\delta_n)t}-e^{(-\mu-b_n)t})}.
\end{equation}
Thus
\begin{align}
Y_n(0)&=\frac{C_n}{b_n+\delta_n}\left[\delta_nF(-\mu+\delta_n)+b_nF(-\mu-b_n)\right],\label{eq:Ymodal}\\
P_n(0)&=\frac{C_n\gamma}{b_n+\delta_n}\left[F(-\mu+\delta_n)-F(-\mu-b_n)\right].\label{eq:Pmodal}
\end{align}
If \(f(t)=e^{\mu t}t v(t)\), the zero-temperature condition is precisely
\begin{equation}\label{eq:paired}
\int_0^T t v(t)\left(e^{-b_nt}+\frac d{b_n^2}e^{\delta_nt}\right)dt=0,\qquad n\ge1.
\end{equation}
The coefficient of the slow term has a plus sign and is \(d/b_n^2\). Its values are tied to the same \(v\) as the fast values.

\section{Bessel-kernel operator and functional setting}\label{sec:setting}
From here the theorem is abstract: fix \(T,d,\mu>0\) and an increasing positive sequence satisfying
\begin{equation}\label{eq:growth}
b_n\asymp n^2,\qquad \inf_n(\sqrt{b_{n+1}}-\sqrt{b_n})>0.
\end{equation}
The gap is a separate hypothesis: growth order alone permits clustered pairs. All spaces are real. Put \(H=L^2((0,T),t\,dt)\) and
\begin{equation}\label{eq:kernel}
H_d(t,s)=\sum_{j=0}^\infty\frac{(dts)^j}{j!(j+1)!},\qquad
(\mathsf Tv)(t)=d\int_0^T H_d(t,s)v(s)s\,ds.
\end{equation}
\begin{lemma}\label{lem:positive}
The operator \(\mathsf T\) is compact, self-adjoint and nonnegative on \(H\). Its quadratic form is
\begin{equation}\label{eq:positivity}
\langle\mathsf Tv,v\rangle_H=d\sum_{j=0}^\infty\frac{d^j}{j!(j+1)!}\left(\int_0^T s^{j+1}v(s)ds\right)^2.
\end{equation}
The inverse \(Q=(I+\mathsf T)^{-1}\) exists with \(\|Q\|\le1\). For \(v\in H\), \(\mathsf Tv\) extends to an entire function of \(t\); consequently \(Qk\) is smooth on \([0,T]\) for smooth \(k\).
\end{lemma}
\begin{proof}
The series converges uniformly on compact sets in \((t,s)\), since \(j!(j+1)!\ge(j!)^2\). Thus the symmetric kernel is bounded on \([0,T]^2\) and belongs to \(L^2((0,T)^2,ts\,dt\,ds)\), proving Hilbert--Schmidt compactness and self-adjointness. Uniform convergence and Cauchy--Schwarz justify integration termwise, giving \eqref{eq:positivity}. Since \(I+\mathsf T\) is bounded below by one and self-adjoint, its range is closed and has zero orthogonal complement; it is onto and \(\|Q\|\le1\). Uniform convergence of the integral series on complex compact sets proves the entire extension. Finally \(Qk=k-\mathsf TQk\) almost everywhere, giving the asserted smooth representative.
\end{proof}
For \(k\in H\) define the correction moments
\begin{equation}\label{eq:B}
(\mathsf Bk)_n=\int_0^T te^{-b_nt}k(t)dt+
\int_T^\infty te^{-b_nt}(\mathsf TQk)(t)dt.
\end{equation}
The second integral is finite by the growth bound proved in \cref{lem:tail}. The identity
\begin{equation}\label{eq:laplacekernel}
\int_0^\infty te^{-bt}H_d(t,s)dt=b^{-2}e^{ds/b},\qquad b>0,\quad 0\le s\le T,
\end{equation}
follows by termwise integration: the integral of \(t^{j+1}e^{-bt}\) is \((j+1)!b^{-j-2}\). Nonnegative terms and then Cauchy--Schwarz justify the interchange. If \(v=Qk\), then \(k=v+\mathsf Tv\), and \eqref{eq:laplacekernel} gives
\begin{equation}\label{eq:exactidentity}
(\mathsf Bk)_n=\int_0^T t v(t)\left(e^{-b_nt}+d b_n^{-2}e^{dt/b_n}\right)dt.
\end{equation}
Thus the tail is not an independent unknown: it is determined by \(k\) through \(Q\).

\section{Main theorem}\label{sec:theorem}
For smooth \(k\), use the representative of \(Qk\) in \cref{lem:positive}, and set
\begin{equation}\label{eq:adm}
\mathcal K_{\rm adm}=\{k\in C^\infty[0,T]: k^{(j)}(0)=0\ (j\ge0),\ (Qk)(T)=0\}.
\end{equation}
\begin{theorem}\label{thm:main}
Assume \eqref{eq:growth}. For every open interval \(I\) with \(\overline I\subset(0,T)\), the real vector space
\begin{equation}\label{eq:kernelspace}
\{k\in C_c^\infty(I):\mathsf Bk=0,\ (Qk)(T)=0\}
\end{equation}
is infinite-dimensional. In particular, \(\ker(\mathsf B|_{\mathcal K_{\rm adm}})\) is infinite-dimensional.
\end{theorem}
The space includes zero; the assertion provides infinitely many linearly independent nonzero corrections. Its proof is given in \cref{sec:biorth,sec:fredholm,sec:endpoint}. It does not assert a numerical formula or control cost.

\section{Biorthogonal construction and smooth right inverse}\label{sec:biorth}
Fix \(I\Subset(0,T)\). Choose \(a,\ell,\varepsilon,\eta>0\) such that \([a,a+\ell+\varepsilon]\subset I\) and \(a+\varepsilon<T-\eta\). Let \(X=L^2(I,t\,dt)\), embedded in \(H\) by zero extension, and define
\begin{equation}\label{eq:EY}
(Ex)_n=\int_I te^{-b_nt}x(t)dt,\qquad
Y=\left\{y:\sum_n e^{2(T-\eta)b_n}|y_n|^2<\infty\right\}.
\end{equation}
The space \(Y\) has its indicated weighted \(\ell^2\) norm. The map \(E:X\to\ell^2\) is bounded: the sum of squares of its coordinate-functional norms is bounded by \(\sum_n\int_I te^{-2b_nt}dt<\infty\).

\begin{proposition}\label{prop:rightinverse}
There is a bounded linear \(S:Y\to X\) whose values have a common compact support in \(I\), belong to \(C_c^\infty(I)\), and satisfy \(ESy=y\) coordinatewise. Its defining series converges in every \(C^j\) norm.
\end{proposition}
\begin{proof}
Let \(g_0>0\) be the infimum in \eqref{eq:growth}, and choose \(0<l\le \min\{\sqrt{b_1},g_0\}\). Apply Theorems 1.2 and 1.5(b) of Fattorini--Russell \cite{FattoriniRussell1974} with \(F(v)=v^2\), \(G(v)=\sqrt v\), and \(\lambda_n=b_n\). The function \(F\) is convex and strictly increasing on \([0,\infty)\), \(F(0)=0\), and \(\int_1^\infty F(v)^{-1}dv=1\). Also \(G(uv)=G(u)G(v)\), while \(\sqrt{b_1}\ge l\) and \(\sqrt{b_{n+1}}-\sqrt{b_n}\ge l\). Thus, for this fixed \(\ell>0\), there are real \(p_n\in L^2(0,\ell)\), extended by zero, satisfying
\begin{equation}\label{eq:pn}
\int_{\mathbb R}e^{-b_mu}p_n(u)du=\delta_{mn},\qquad \|p_n\|_2\le B\exp(K\sqrt{b_n}).
\end{equation}
The constants \(B,K\) depend on the fixed gap and time, but not on \(n\). No upper gap or precise quadratic expansion is required for this source estimate. The subsequent summability uses \(b_n\asymp n^2\). No time reversal is needed here.
Choose a nonzero nonnegative \(\rho\in C_c^\infty(0,\varepsilon)\), and put \(\rho_n=\int e^{-b_nu}\rho(u)du\). Then \(\rho_n\ge e^{-\varepsilon b_n}\int\rho>0\). Define
\begin{equation}\label{eq:psi}
\psi_n(t)=\frac{e^{ab_n}}{t\rho_n}[p_n(\cdot-a)*\rho](t).
\end{equation}
The convolution has a common compact support inside \(I\), where division by \(t\) is smooth. Factoring its Laplace transform gives \(\int_I te^{-b_mt}\psi_n(t)dt=\delta_{mn}\): on the diagonal the factors \(e^{ab_n}\), \(e^{-ab_m}\), \(\rho_m\) and \(\rho_n^{-1}\) cancel; off the diagonal \eqref{eq:pn} vanishes. By Cauchy--Schwarz on each convolution derivative and Leibniz' rule, for every \(j\ge0\),
\begin{equation}\label{eq:psibound}
\|\psi_n\|_{C^j}\le C_j e^{(a+\varepsilon)b_n+C\sqrt{b_n}}.
\end{equation}
Define \(Sy=\sum_n y_n\psi_n\). With \(c=T-\eta-a-\varepsilon>0\), Cauchy--Schwarz bounds the sum of \(C^j\) norms by \(C_j\|y\|_Y(\sum_n e^{-2cb_n+2C\sqrt{b_n}})^{1/2}\), which is finite using \(b_n\asymp n^2\). Hence the derivative series is uniformly absolutely convergent for every \(j\), with common compact support. Termwise integration proves \(ESy=y\). The right inverse is only asserted on \(Y\), not on all of \(\ell^2\).
\end{proof}

\section{Compact tail and Fredholm intersection}\label{sec:fredholm}
\begin{lemma}\label{lem:tail}
Define \(R:X\to Y\) by
\begin{equation}\label{eq:R}
(Rx)_n=\int_T^\infty te^{-b_nt}(\mathsf TQx)(t)dt.
\end{equation}
This map is Hilbert--Schmidt. Hence \(SR:X\to X\) is compact.
\end{lemma}
\begin{proof}
For \(z\ge0\), \(\sum_{j\ge0}z^j/(j!)^2\le e^{2\sqrt z}\), since the sum is dominated by the square of \(\sum_j(\sqrt z)^j/j!\). Thus, for \(t\ge T\),
\begin{equation}\label{eq:tailgrowth}
|(\mathsf TQx)(t)|\le d\|Qx\|_H\left(\int_0^T s|H_d(t,s)|^2ds\right)^{1/2}
\le C e^{2\sqrt{dTt}}\|x\|_X.
\end{equation}
Writing \(t=T+u\), and using \(b_n\ge b_1>0\),
\[
\int_T^\infty te^{-b_nt}e^{2\sqrt{dTt}}dt
\le e^{-Tb_n}\int_0^\infty(T+u)e^{-b_1u+2\sqrt{dT(T+u)}}du
\le C e^{-Tb_n}.
\]
The final integral is finite because its exponent is \(-b_1u+O(\sqrt u)\). The weighted coordinate functionals of \(R\) therefore have squared norms bounded by \(C^2e^{-2\eta b_n}\), whose sum is finite. For a map from a Hilbert space into a weighted \(\ell^2\) space this is exactly the Hilbert--Schmidt criterion. Composition with bounded \(S\) is compact.
\end{proof}
\begin{proposition}\label{prop:fredholm}
There is an infinite-dimensional subspace of \(C_c^\infty(I)\) on which \(E+R=0\) coordinatewise.
\end{proposition}
\begin{proof}
Let \(A=I_X+SR\). The Fredholm alternative \cite{Howard2020} gives closed range of finite codimension; neither invertibility of \(A\) nor smallness of \(SR\) is needed. We first show that
\[
N=\{h\in C_c^\infty(I):Eh=0\}
\]
is infinite-dimensional. Choose a closed interval \(J\) in the interior of \(I\). By \(b_n\asymp n^2\), \(\sum_n b_n^{-1}<\infty\). The full M\"untz theorem away from zero, \cite[Theorem~29]{Almira2007}, applied after \(z=e^{-t}\), shows that the fast exponentials are not dense in \(C(J)\). Hahn--Banach and the Riesz representation theorem \cite{RieszNotes} provide a nonzero finite signed measure \(\nu\) on \(J\) annihilating all of them.

Choose a nonzero compactly supported smooth bump \(\rho\) so small that \(J+\operatorname{supp}\rho\Subset I\). The function \(w=\nu*\rho\) is smooth and compactly supported, and its Laplace moments factor into those of \(\nu\) and \(\rho\), hence vanish. It is nonzero: the Fourier transforms of both compactly supported factors are nonzero entire functions, and their product is not identically zero. Then \(h=w/t\) is smooth, compactly supported and belongs to \(N\). Repeating the construction in pairwise disjoint subintervals of \(I\) gives linearly independent \(h\)'s.

The quotient \(X/\operatorname{ran}A\) is finite-dimensional. Therefore the kernel of the quotient map restricted to \(N\), namely \(N\cap\operatorname{ran}A\), is infinite-dimensional. Choose linearly independent \(h_i\) there and solve \(Ax_i=h_i\) in \(X\). Since \(x_i=h_i-SRx_i\) as an \(X\)-identity, \(x_i\) has a \(C_c^\infty(I)\) representative. Moreover, \(Ex_i=Eh_i-ESRx_i=-Rx_i\), because \(Rx_i\in Y\) and \(ES=I_Y\). The \(x_i\) are independent: applying \(A\) sends any linear dependence to the corresponding dependence among the \(h_i\). Their span gives the proposition. No finite-mode approximation is used.
\end{proof}

\section{Endpoint constraint and proof of the theorem}\label{sec:endpoint}
\begin{lemma}\label{lem:endpoint}
The functional
\begin{equation}\label{eq:endpoint}
L(x)=-d\int_0^T H_d(T,s)(Qx)(s)s\,ds,\qquad x\in X,
\end{equation}
is bounded. If \(x\in C_c^\infty(I)\), then \(L(x)=(Qx)(T)\), with the smooth representative.
\end{lemma}
\begin{proof}
The kernel row \(dH_d(T,\cdot)\) belongs to \(H\) and \(Q\) is bounded, so Cauchy--Schwarz gives \(|L(x)|\le C\|x\|_X\). Point evaluation is not defined for an arbitrary element of \(H\). However, the zero extension of \(x\in X\) vanishes almost everywhere near \(T\), and \(Qx=x-\mathsf TQx\) has the canonical representative \(-\mathsf TQx\) there. For smooth compactly supported \(x\), this agrees with the smooth representative on all of \([0,T]\), giving the claimed endpoint value.
\end{proof}
\begin{proof}[Proof of \cref{thm:main}]
Let \(V\subset C_c^\infty(I)\) be the infinite-dimensional space from \cref{prop:fredholm}. For every \(k\in V\), \((E+R)k=\mathsf Bk=0\), by the zero extension in \eqref{eq:B}. The restriction of one real linear functional \(L\) to \(V\) has a kernel of codimension at most one. Hence \(V\cap\ker L\) is infinite-dimensional, and \cref{lem:endpoint} gives \((Qk)(T)=0\) for each of its members. Such \(k\) also vanishes near zero with all derivatives, so it lies in \(\mathcal K_{\rm adm}\). This proves the theorem.
\end{proof}

\section{Conditional terminal-state interpretation}\label{sec:interpretation}
Take \(b_n\) from \eqref{eq:roots}, and a nonzero \(k\) in \eqref{eq:kernelspace}. Put \(v=Qk\), \(f(t)=e^{\mu t}tv(t)\), \(a(s)=f(-s)\). Then \(v\) is smooth, \(f(0)=0\), and \(f(T)=0\) by the endpoint condition. Under the smooth-lift solution framework in \cref{sec:model}, \eqref{eq:exactidentity} and \eqref{eq:Ymodal} imply
\begin{equation}\label{eq:temperature}
Y_n(0)=\frac{C_nb_n}{b_n+\delta_n}(\mathsf Bk)_n=0.
\end{equation}
All coefficients in the odd basis vanish, so the odd terminal temperature is zero in \(L^2(0,L)\). The zero-temperature relation in \eqref{eq:Ymodal} gives \(F(-\mu-b_n)=-d b_n^{-2}F(-\mu+\delta_n)\). Substitution in \eqref{eq:Pmodal} yields
\begin{equation}\label{eq:memory}
P_n(0)=\frac{C_n\gamma}{b_n}F(-\mu+\delta_n)
=\frac{C_n\gamma}{b_n}\int_0^T te^{\delta_nt}v(t)dt.
\end{equation}
No terminal condition \(P(0)=0\) is imposed.

Because \(Q\) is injective, \(v\ne0\). Each weight in \eqref{eq:paired} is positive on \((0,T)\), so the real continuous functions \(v\) and \(f\) must change sign. Since \(k\) vanishes near both endpoints, \(v=-\mathsf Tv\) there, and \(f\) is analytic near each endpoint. If \(f\) vanished on \((0,\epsilon)\), then \(\mathsf Tv\) would vanish there and hence everywhere by analyticity. By \eqref{eq:positivity}, all weighted moments of \(v\) would vanish; polynomial density in \(H\) would imply \(v=0\), a contradiction. Thus, in physical time, the control has no terminal switch-off gap.

The entire function \(G(z)=\int_0^T te^{zt}v(t)dt\) is nonzero: otherwise all weighted polynomial moments vanish and \(v=0\). Since \(\delta_n\to0\), its zeros cannot contain infinitely many \(\delta_n\) unless \(G\equiv0\). From \eqref{eq:memory}, \(C_n\ne0\) and \(\gamma\ne0\), so all but finitely many terminal memory coefficients are nonzero. Finally, \(f\) cannot be real analytic on a connected neighborhood of \([0,T]\): the quotient \(e^{-\mu t}f(t)/t\) has an analytic removable singularity at zero, making \(k=v+\mathsf Tv\) analytic and identically zero because it vanishes near zero. These are assertions about this linear odd solution class only.

\section{Terminal memory: necessary trace conditions}\label{sec:trace}
For a compatible smooth boundary input, let \(f(t)=a(-t)\), and suppose \(Y(\cdot,0)=0\). Put \(c_n=C_n\gamma/b_n\ne0\), \(\delta_n=d/b_n\), \(p_n=P_n(0)\), and
\[
G(z)=\int_0^T e^{zt}e^{-\mu t}f(t)\,dt.
\]
For the constructed controls this agrees with the function in \cref{sec:interpretation}; the modal identity \eqref{eq:memory} gives \(p_n=c_nG(\delta_n)\). The function \(G\) is entire of exponential type at most \(T\). Set \(r_n=p_n/c_n\).

\begin{proposition}[Implicit exact characterization]\label{prop:implicit}
In the smooth compatible class specified in \cref{sec:model}, a sequence \(p\in\ell^2\) is a terminal memory sequence of a zero-start control with zero terminal temperature if and only if there exists \(h\in C^\infty([0,T])\), \(h(0)=h(T)=0\), such that, with \(G(z)=\int_0^T e^{zt}h(t)dt\),
\[
G(-b_n)+\frac{d}{b_n^2}G(\delta_n)=0,\qquad
p_n=c_nG(\delta_n),\quad n\ge1.
\]
This is an exact implicit moment description, not an intrinsic sequence-space characterization or a statement about closure.
\end{proposition}
\begin{proof}
For any such control take \(h=e^{-\mu t}f(t)\). The temperature identity \eqref{eq:Ymodal} gives the first equation and \eqref{eq:memory} gives the second. Conversely set \(f=e^{\mu t}h\) and \(a(s)=f(-s)\). It is smooth and endpoint-compatible. The two projector identities \eqref{eq:Ymodal}--\eqref{eq:Pmodal} imply that every temperature coefficient vanishes and every memory coefficient is \(p_n\). Completeness of the odd basis and the smooth-lift solution framework give the terminal fields in \(L^2\).
\end{proof}

\begin{proposition}[Accumulating-zero mechanism]\label{prop:fullzero}
For the frozen zero-start model with the stated smooth compatible boundary input, if \(Y(\cdot,0)=P(\cdot,0)=0\), then \(a=0\). In particular, the memory of any nonzero constructed temperature-null input is nonzero. The use of zeros accumulating at a finite point is the mechanism of \cite[Obstruction 1, p.~6]{IvanovPandolfi2009}, not a novelty claim.
\end{proposition}
\begin{proof}
The terminal identities give \(G(\delta_n)=0\) for all \(n\). As the distinct \(\delta_n\) converge to zero, the identity theorem gives \(G\equiv0\). Hence all polynomial moments of \(e^{-\mu t}f(t)\) vanish. Polynomial density in \(L^2(0,T)\) yields \(f=0\), and therefore \(a=0\).
\end{proof}
\begin{proposition}[Necessary asymptotics and a norm lower bound]\label{prop:asymptotic}
For every integer \(m\ge0\), put \(g_j=\int_0^T t^j e^{-\mu t} f(t)dt\). Then
\[
r_n=\sum_{j=0}^{m}\frac{g_j}{j!}\delta_n^j+O(\delta_n^{m+1}),\qquad
|g_j|\le\left(\int_0^T t^{2j}e^{-2\mu t}dt\right)^{1/2}\|a\|_{L^2(-T,0)}.
\]
In particular \(r_n\to g_0\) and every smooth compatible input producing a prescribed zero-temperature memory sequence satisfies
\[
\|a\|_{L^2(-T,0)}\ge \frac{|\lim_n p_n/c_n|}{\left(\int_0^T e^{-2\mu t}dt\right)^{1/2}}.
\]
The normalized sequences determine \(G\) and hence the input uniquely whenever a solution exists.
\end{proposition}
\begin{proof}
Taylor expansion of the entire integral and boundedness of \(t\in[0,T]\) give the expansion, with remainder controlled by \(\|f\|_2\). Cauchy--Schwarz gives the coefficient bounds. If two functions agree at all \(\delta_n\), their difference vanishes identically by the identity theorem, so polynomial density proves equality of the inputs. The time reversal preserves the \(L^2\) norm.
\end{proof}
Since \(b_n\asymp n^2\) and \(C_n\asymp n\), the sequence \((c_n)\in\ell^2\). Thus \(p_n=c_n(-1)^n\) is an \(\ell^2\) memory sequence violating the necessary convergence condition and is not exactly reachable with zero temperature in this class. This gives no conclusion about its membership in the \emph{closure} of the range. An intrinsic description of the exact range, its density in a specified weighted space and sharp control costs remain open; the trace conditions alone are not sufficient. The coefficient bounds are elementary and no novelty is asserted.

\begin{corollary}[Elementary instability, conditional on the kernel theorem]
There are smooth compatible real boundary controls of unit \(L^2(-T,0)\) norm, each giving zero terminal temperature, whose terminal memory norms tend to zero in \(L^2(0,L)\). This is not a nonzero control with exactly zero terminal full state and is not a control-cost bound for prescribed states.
\end{corollary}
\begin{proof}
Fix \(I\Subset(0,T)\), put \(X=L^2(I,t\,dt)\), and let \(V\subset C_c^\infty(I)\) be the infinite-dimensional correction space furnished by \cref{thm:main}. For \(x\in V\), set \(v=Qx\), \(f(t)=e^{\mu t}tv(t)\). The memory map \(x\mapsto(c_n\int_0^T te^{\delta_nt}(Qx)(t)dt)_n\) is compact from \(X\) to \(\ell^2\): its Taylor truncations in \(\delta_n\) have finite rank, and converge in operator norm since \((c_n)\in\ell^2\), \(0<\delta_n\le d/b_1\), and \(Q\) is bounded on \(H\).

On the other hand, with \(\alpha=\inf I>0\), \(\|v\|_{X}\le\alpha^{-1/2}\|f\|_2\). The Bessel kernel is bounded on \([0,T]^2\), and \(sv(s)=e^{-\mu s}f(s)\), so \(\|\mathsf Tv\|_X\le C\|f\|_2\). Since \(x=v+\mathsf Tv\) on \(I\), \(\|x\|_X\le C_I\|f\|_2\). Apply Gram--Schmidt to independent elements of \(V\), obtaining an \(X\)-orthonormal sequence \(x_j\in V\). Compactness makes the corresponding memory coefficients tend to zero, while \(\|f_j\|_2\ge C_I^{-1}\). Normalize by \(\|f_j\|_2\), reverse time and use the modal reconstruction. Constants depend on the fixed model and interval.
\end{proof}

\section{Discussion and limitations}\label{sec:discussion}
The construction combines two tools that address different parts of the problem: M\"untz incompleteness supplies fast-null seeds, while quantitative biorthogonality solves fast interpolation for the tail's rapidly decaying coordinates. The Fredholm step is applied to \(I+SR\) on the interior-support space, not to the original coupled moment map. Its range is used exactly, not replaced by its closure. The construction is existential and gives no uniform quantitative control cost or prescribed Gevrey order.

The asymptotic condition and the uniform square-root gap are both stated because the former alone does not imply the latter. The precise external estimate and its specialization require independent checking; our use of it is not a new proof of the cited result. The PDE interpretation additionally requires the specified frozen coefficients, zero initial odd fields, compatible boundary values and the smooth-lift solution class. It is not a nonlinear state-preparation theorem, a controllability-to-rest theorem, a proof of physical reachability, or an inverse result for the full system. The literature review is not exhaustive and establishes no chronology or novelty. The proof and model interpretation have not been independently certified.

\section{Conclusion}
Under quadratic spectral growth and a separate positive square-root gap, the dependent-tail representation and a weighted smooth right inverse yield an infinite-dimensional kernel of the exact coupled moment map, even after the endpoint constraint. In the stated frozen linear odd model these inputs give zero terminal temperature and generally nonzero terminal memory. Independent mathematical, model and literature review remain necessary before any submission or priority claim.

\bibliographystyle{unsrt}
\bibliography{references}
\end{document}
